% ==================== 章节1：深夜小茶会 ====================
\chapter{深夜小茶会}

% --- 小剧场 ---
\section*{初星学园——11:40P.M. 活动室}
\phantomsection
\addcontentsline{toc}{section}{小剧场}

\skitillustration{figures/skits/ch01-night-tea.png}

深夜的活动室只剩一盏灯。窗外的初星学园安静得像被谁按下暂停键。

Producer 在桌边收训练记录，笔尖划过纸面。広抱着热茶，坐在旁边，看起来像快睡着了，又像正在读取什么很长的日志。

P：今天先到这里。你已经很累了。

広：嗯。身体的电量，确实很低。

P：你还知道。

広：知道。不代表会停止观察。

Producer 叹了口气，把最后一页记录夹进文件夹。

P：如果有个小助手就好了。至少能在你倒下以后，把后面的事接着做完。

広抬起眼。

P：我只是随口说。

広：随口说出来的愿望，也可以成为需求。

第二天，広把一张手写纸放到桌上。标题写得很认真：

\begin{center}
小広计划：第一版需求草案
\end{center}

下面分着“目标”“限制”“失败条件”“暂不考虑事项”。字很小，结构却清楚得像正式项目文档。

P：你真的写了？

広：嗯。你昨天说出口了。

P：这里为什么有一条“本体容易倒下”？

広：系统风险。

P：这行划掉。

広把纸往自己这边轻轻拉了一点。

広：不可以。风险不能因为看起来可怜就消失。

P：你把自己写得像硬件缺陷。

広：硬件缺陷也可以设计边界。这样比较温柔。

Producer 没有立刻反驳。他看着那行字，又看向她手里的热茶。

P：为什么来真的？

広：因为你说出口了。……而且很不适合我。

她停了一下，像是在确认结论。

広：这样很好。

\clearpage

% 正文内容
\section{深夜小茶会}

这本小册子的起点，不是实验室，也不是发布会。

它从一杯快凉掉的热茶开始，从一句本来不该被认真对待的话开始。Producer 以为自己只是在担心広的身体，広却把那句话翻译成一门课题：如果想做一个能把事情继续接下去的“小広”，要先理解 AI 是怎样从文字、计划、身体和失败里一点点长出来的。

所以，后面的每一章都不是单纯的技术讲义。

它们更像放学后的补课记录：有白板，练习室，门把手，胶带，也有一台很慢、很笨、终于把世界往前推了一点的试作机。
